\section{Install the tools}
\label{sec:tools}

\subsection{What is installed, and why}

Four programs are installed in this section (Table~\ref{tab:tools}). Three of them are installed
with \term{winget}, the package manager built into Windows, which downloads a program from its
publisher and installs it in response to a single command. The fourth, MiKTeX, is installed with
its own installer because two of its options must be chosen by hand.

\begin{table}[H]
\centering
\begin{tabular}{C{2.0cm}L{6.4cm}C{2.9cm}}
\toprule
\hd{Program} & \multicolumn{1}{c}{\hd{Purpose in this project}} & \hd{Commands it provides} \\
\midrule
Git & Downloads the repository and exchanges changes with GitHub & \cmd{git} \\
Python 3.12 & Runs every program of the pipeline & \cmd{py}, \cmd{python}, \cmd{pip} \\
MiKTeX & Turns \LaTeX{} text into a PDF; supplies the fonts of the printed material &
  \cmd{pdflatex} \\
Poppler & Turns a PDF page into a picture, for the simulated photographs & \cmd{pdftoppm} \\
\bottomrule
\end{tabular}
\caption{The tools installed in this section.}
\label{tab:tools}
\end{table}

\begin{note}[Why Python 3.12 and not the newest version]
The neural-network library TensorFlow is published for particular Python versions only. The
project was developed and tested with Python 3.12 and TensorFlow 2.21. Installing 3.12 does not
disturb any other Python version that may already be on the laptop.
\end{note}

\subsection{Check the package manager}

\step{Confirm that winget is present}
Open a PowerShell window (Step~\ref{st:openps}) and run:
\begin{ps}
winget --version
\end{ps}
\begin{out}
v1.11.430
\end{out}
\checkpoint{A version number beginning with \texttt{v1} is printed. The digits after it do not
matter.}

\begin{iffail}[If PowerShell answers that ``winget'' is not recognized]
Click \ui{Start}, type \code|store| and open \ui{Microsoft Store}. Search for
\ui{App Installer}, open its page, and click \ui{Get} or \ui{Update}. When the Store has
finished, close the PowerShell window, open a new one, and repeat this step.
\end{iffail}

\subsection{Git}

\step{Install Git}
\begin{ps}
winget install --id Git.Git -e --source winget
\end{ps}
The first time winget is used it may ask whether you agree to the source agreement terms; type
\code|Y| and press \key{Enter}. A few seconds later Windows shows a dialog titled \ui{User
Account Control} that asks whether the app may make changes to the device; click \ui{Yes}. The
installer then runs by itself for one or two minutes.
\begin{out}[What you should see at the end]
Successfully installed
\end{out}
\checkpoint{The last line is \emph{Successfully installed} and a new prompt has appeared.}

\subsection{Python}

\step{Install Python 3.12}
\begin{ps}
winget install --id Python.Python.3.12 -e --source winget
\end{ps}
If the \ui{User Account Control} dialog appears, click \ui{Yes}.
\checkpoint{The last line is \emph{Successfully installed}.}

\subsection{MiKTeX}

\step{Download the MiKTeX installer}
In the browser, open
\begin{center}\url{https://miktex.org/download}\end{center}
Make sure the tab \ui{Windows} is selected, and under the heading \ui{Installer} click
\ui{Download}. The file is named \file{basic-miktex-} followed by a version number and
\file{-x64.exe}, and is about 140\,MB.
\checkpoint{The file is in your \file{Downloads} folder.}

\step{Run the installer}
Open the \file{Downloads} folder (in File Explorer, \key{Windows}+\key{E}, then
\ui{Downloads}) and double-click the downloaded file. The installer is a sequence of pages;
answer them as in Table~\ref{tab:miktex} and click \ui{Next} after each.

\begin{table}[H]
\centering
\begin{tabular}{C{3.2cm}L{8.4cm}}
\toprule
\hd{Page} & \multicolumn{1}{c}{\hd{What to choose}} \\
\midrule
Copying conditions & Tick \emph{I accept the MiKTeX copying conditions} \\
Installation scope & Choose the \emph{private (per-user)} installation, which MiKTeX itself
  recommends and which needs no administrator rights \\
Installation directory & Leave the proposed folder unchanged \\
Settings & Preferred paper: \emph{Letter}. Install missing packages on-the-fly:
  \emph{Always} \\
Review & Click \ui{Start}; the installation takes a few minutes \\
Finish & Click \ui{Close} \\
\bottomrule
\end{tabular}
\caption{Answers for the pages of the Basic MiKTeX Installer.}
\label{tab:miktex}
\end{table}

\rationale{The pipeline runs \cmd{pdflatex} in the background. With the setting \emph{Always},
MiKTeX fetches any missing component silently. With the default setting it would instead open a
dialog and wait for an answer that the pipeline cannot give, and the operation would time out.}

\checkpoint{The installer has closed, and the Start menu contains an entry \emph{MiKTeX
Console}.}

\step{Let MiKTeX update itself once}
Click \ui{Start}, type \code|miktex console| and open \ui{MiKTeX Console}. On its first page
click \ui{Check for updates}. If updates are found, click \ui{Update now} and wait until the
progress display has finished; then close MiKTeX Console.
\checkpoint{MiKTeX Console reports that the installation is up to date (or you have installed
the updates it offered).}

\subsection{Poppler}

\step{Install Poppler}
Return to the PowerShell window and run:
\begin{ps}
winget install --id oschwartz10612.Poppler -e --source winget
\end{ps}
\checkpoint{The last line is \emph{Successfully installed}.}

\subsection{Verify all four tools}

\step{Open a new PowerShell window}\label{st:newafterinstall}
Close the PowerShell window (type \code|exit| and press \key{Enter}), then open a new one as in
Step~\ref{st:openps}.

\rationale{A PowerShell window learns the list of installed programs at the moment it opens.
The window used for the installations therefore does not yet know the new commands; a new
window does.}

\step{Make MiKTeX install missing components without asking}\label{st:autoinstall}
\begin{ps}
initexmf "--set-config-value=[MPM]AutoInstall=1"
\end{ps}
\checkpoint{The command prints nothing and a new prompt appears. (It repeats the choice
\emph{Always} of Table~\ref{tab:miktex}, so that the setting is right even if the installer page
was answered differently.)}

\step{Check the four commands}\label{st:verifytools}
Run the following four commands one after the other. Each prints a line that names the program
and its version; only the first line of each answer matters, and your version numbers may be
higher than the ones shown.

\begin{ps}
git --version
\end{ps}
\begin{out}
git version 2.52.0.windows.1
\end{out}

\begin{ps}
py -3.12 --version
\end{ps}
\begin{out}
Python 3.12.10
\end{out}

\begin{ps}
pdflatex --version
\end{ps}
\begin{out}[What you should see (first line)]
MiKTeX-pdfTeX 4.21 (MiKTeX 25.12)
\end{out}

\begin{ps}
pdftoppm -v
\end{ps}
\begin{out}[What you should see (first line)]
pdftoppm version 25.07.0
\end{out}

\checkpoint{All four commands answered with a version. The Python version begins with
\texttt{3.12}.}

\begin{iffail}[If one of the four commands is ``not recognized'']
Close the window, open a new one and try that command once more. If it still fails, the program
was not installed or Windows cannot find it: repeat the installation step of that program and
watch for an error message. For \cmd{pdftoppm} in particular, Section~\ref{sec:trouble} gives
three commands that repair the search path.
\end{iffail}
